\input{../header}
\usetikzlibrary{shapes.geometric}

\begin{document}
%


%\onehalfspacing
\allowdisplaybreaks
%##################################################################
\section{Warmup: Similar triangles}

One fun trick that we often use to relate variables in related rates problems is a thing you probably learned in high school but maybe forgot about: \textit{similar triangles}. This warmup will help you see what this is and why it works.

\begin{enumerate}[leftmargin=15pt]
    \item Draw a big triangle, any shape you want. Could be equilateral, could be right, could be scalene. Label its sides $A$, $B$, and $C$, and measure each side with a ruler.
    \vfill
    \item Inside the big triangle, draw a line parallel to one of the sides. You have thus made a smaller triangle inside the big triangle. Label the sides of this triangle $a$, $b$, and $c$, corresponding to your labeling of the big triangle, and measure each side with a ruler.
    \item Compute a bunch of ratios, then make some observations:
    \begin{align*}
        \dfrac{A}{a} &= & \dfrac{B}{b} &= & \dfrac{C}{c} &= \\\\
        \dfrac{A}{B} &= & \dfrac{A}{C} &= & \dfrac{B}{C} &= \\\\
        \dfrac{a}{b} &= & \dfrac{a}{c} &= & \dfrac{b}{c} &=
    \end{align*}
\end{enumerate}


\end{document}